 \chapter{Ergebnisse}\label{ch:ergebnisse}

In diesem Kapitel werden die Resultate der durchgeführten Experimente, Berechnungen oder Analysen präsentiert.
Wichtig ist eine klare, strukturierte und objektive Darstellung.

\begin{itemize}
\item Beschreiben Sie zunächst den Versuchsaufbau und die Rahmenbedingungen.
\item Führen Sie Vergleichs- oder Referenzergebnisse an (Baseline, theoretische Erwartung, Stand der Technik).
\item Stellen Sie quantitative Resultate in Tabellen oder Diagrammen dar. Nutzen Sie aussagekräftige Beschriftungen und Einheiten.
\item Diskutieren Sie kurz die wichtigsten Beobachtungen (ohne Interpretationen – die folgen im nächsten Kapitel).
\end{itemize}

Bei empirischen Arbeiten sind Mittelwerte, Standardabweichungen oder Konfidenzintervalle zu nennen; bei Simulationsstudien Laufzeiten, Genauigkeiten oder Fehlermaße.


